\documentclass[12pt, a4paper]{article}

\usepackage[utf8]{inputenc}
\usepackage[T1]{fontenc}
\usepackage[ngerman]{babel}
\usepackage{amsmath}
\usepackage{geometry}
\usepackage{parskip}

\geometry{top=2.5cm, bottom=2.5cm, left=2.8cm, right=2.8cm}

\begin{document}

% ---------- Header ----------
\begin{minipage}[t]{0.55\textwidth}
  Institut für Computerphysik\\
  Universität Stuttgart
\end{minipage}%
\hfill
\begin{minipage}[t]{0.4\textwidth}
  \raggedleft
  Physik auf dem Computer\\
  SS 2026
\end{minipage}

\rule{\textwidth}{0.4pt}

\vspace{0.6em}

{\LARGE\bfseries Übungsblatt 2: Perkolation und Clustereigenschaften}

\vspace{1em}

% ---------- Intro ----------
In diesem Übungsblatt implementieren Sie zunächst den Hoshen-Kopelman-Algorithmus
zur Cluster-Identifizierung und untersuchen anschließend drei charakteristische
Größen des Phasenübergangs am kritischen Punkt numerisch.

Für ein zweidimensionales Quadratgitter liegt der kritische Punkt bei $p_c \approx 0.5927$.


\vspace{0.5em}

% ---------- Tasks ----------
\begin{enumerate}

% ---- 1 ----
\item Implementieren Sie den \textbf{Hoshen-Kopelman-Algorithmus} zur Cluster"=Beschriftung. Implementieren Sie dabei \textbf{UnionFind} selbst.

% ---- 2 ----
\item Berechnen Sie die \textbf{Clustergrößenverteilung}
\[
  n_s(p) \;=\; \frac{\text{Anzahl Cluster der Größe }s}{L^2}
\]
für $p = p_c \approx 0.5927$ auf einem $200\times 200$-Gitter.
Mitteln Sie über mindestens 200 unabhängige Konfigurationen.

Tipp: Perkolierendes Cluster ausschließen!

Stellen Sie $n_s$ in einem doppelt-logarithmischen Diagramm dar.
Am kritischen Punkt gilt das Potenzgesetz
\[
  n_s(p_c) \;\propto\; s^{-\tau}, \qquad \tau = \tfrac{187}{91} \approx 2.055\,,
\]
mit einem universellen Exponenten, der nur von der Raumdimension abhängt.
Ergänzen Sie Ihren Plot um diese Referenzgerade.
Können Sie den Zusammenhang bestätigen?

Wiederholen Sie die Rechnung für $L = 50,\,100,\,200$ in einem gemeinsamen Plot.
Was beobachten Sie und warum?

% ---- 3 ----
\item Berechnen Sie die \textbf{mittlere Clustergröße}
\[
  S(p) \;=\; \frac{\displaystyle\sum_s s^2\, n_s(p)}{\displaystyle\sum_s s\, n_s(p)},
\]
wobei die Summe nur über \emph{nicht}-perkolierende Cluster läuft.
$S(p)$ divergiert am kritischen Punkt, da die typische Clustergröße dort unbegrenzt
anwächst.

Plotten Sie $S(p)$ als Funktion von $p \in [0.3,\,0.75]$
für die Gittergrößen $L = 32,\,64,\,128$.
Mitteln Sie pro $p$-Wert über mindestens 100 Konfigurationen.

Schätzen Sie $p_c$ aus dem Plot und vergleichen Sie mit dem theoretischen Wert
$p_c \approx 0.5927$.

% % ---- 4 ----
% \item Die \textbf{Dichte des perkolierenden Clusters}
% \[
%   P(p) \;=\; \frac{\text{Größe des perkolierenden Clusters}}{\text{Anzahl besetzter Plätze}}
% \]
% ist der Ordnungsparameter des Phasenübergangs: $P = 0$ für $p < p_c$ und $P > 0$
% für $p > p_c$.
% Falls kein perkolierender Cluster existiert, setzen Sie $P = 0$.

% Plotten Sie $P(p)$ für dieselben Gittergrößen und $p$-Werte wie in Aufgabe~3.

% % ---- 5 ----
% \item[\textbf{5.}] \textbf{Bestimmung von $p_c$}

% Der kritische Punkt lässt sich aus den Plots der Aufgaben 3 und 4 schätzen.
% ggf. als Hinweis:
% \begin{itemize}
%   \item \textit{Aus $S(p)$:} Das Maximum von $S(p)$ verschiebt sich mit wachsendem $L$
%         gegen $p_c$.
%   \item \textit{Aus $P(p)$:} Die Kurven für verschiedene $L$ schneiden sich näherungsweise
%         bei $p_c$, da $P(p_c, L)$ schwach $L$-abhängig ist.
% \end{itemize}
% Schätzen Sie $p_c$ aus beiden Plots und vergleichen Sie mit dem theoretischen Wert
% $p_c \approx 0.5927$.
% Welche Methode ist genauer, und warum?

\end{enumerate}

\end{document}
